\chapter{Marco teórico y estado del arte}

\section{Mercados de activos digitales}

% ponytail: pendiente de desarrollar con referencias y ejemplos concretos del dominio CS2.

\subsection{Activos virtuales en Counter-Strike 2}

% ponytail: explicar skins, cajas, rareza, estado, liquidez y mercados externos.

\subsection{Fricciones de mercado}

% ponytail: comisiones, divisas, retrasos, spread, volumen y riesgo de inventario.

\section{Aprendizaje por refuerzo}

% ponytail: explicar MDP, política, recompensa, retorno y exploración sin convertirlo en manual.

\subsection{Procesos de decisión de Markov}

% ponytail: estados, acciones, transición y recompensa.

\subsection{Política, retorno y exploración}

% ponytail: conectar aprendizaje con decisiones secuenciales de cartera.

\section{Aprendizaje por refuerzo multiagente}

% ponytail: introducir MARL, coordinación, cooperación y descentralización.

\subsection{Cooperación y coordinación entre agentes}

% ponytail: explicar por qué dividir el mercado entre agentes puede ayudar.

\subsection{Entrenamiento centralizado y ejecución descentralizada}

% ponytail: incluir CTDE solo si se usa de verdad en el proyecto.

\section{Toma de decisiones en carteras}

% ponytail: conectar rentabilidad, riesgo, liquidez, comisiones y horizonte temporal.

\subsection{Rentabilidad y riesgo}

% ponytail: definir métricas básicas antes de métricas sofisticadas.

\subsection{Liquidez y costes de transacción}

% ponytail: explicar por qué una oportunidad bruta puede no ser rentable.

\section{Aplicaciones relacionadas}

% ponytail: revisar trabajos cercanos solo cuando aporten algo al proyecto.
